%%%%%%%%%%%%%%%%%%%%%%%%%%%%%%%%%%%%%%%%%%%%%%%%%%%%%%%%%%%%%%%%%%%%
%% I, the copyright holder of this work, release this work into the
%% public domain. This applies worldwide. In some countries this may
%% not be legally possible; if so: I grant anyone the right to use
%% this work for any purpose, without any conditions, unless such
%% conditions are required by law.
%%%%%%%%%%%%%%%%%%%%%%%%%%%%%%%%%%%%%%%%%%%%%%%%%%%%%%%%%%%%%%%%%%%%

\documentclass{beamer}
\usetheme[faculty=ped,basePath=../fibeamer,logo=../fibeamer/logo/mu/Logo-UC-3.png]{fibeamer}
\graphicspath{{../resources/}}

\usepackage{algorithm}
\usepackage{algpseudocode}

\usetheme[faculty=ped]{fibeamer}
\usepackage[utf8]{inputenc}
\usepackage[
  main=english, %% By using `czech` or `slovak` as the main locale
                %% instead of `english`, you can typeset the
                %% presentation in either Czech or Slovak,
                %% respectively.
  czech, slovak %% The additional keys allow foreign texts to be
]{babel}        %% typeset as follows:
%%
%%   \begin{otherlanguage}{czech}   ... \end{otherlanguage}
%%   \begin{otherlanguage}{slovak}  ... \end{otherlanguage}
%%
%% These macros specify information about the presentation
\title{Ciclos e iteración (for, while, range)} %% that will be typeset on the
\subtitle{Semestre I-2026} %% title page.
\author{Profesora: Andreina Cota Vidaurre}
%% These additional packages are used within the document:
\usepackage{ragged2e}  % `\justifying` text
\usepackage{booktabs}  % Tables
\usepackage{tabularx}
\usepackage{tikz}      % Diagrams
\usetikzlibrary{calc, shapes, backgrounds}
\usetikzlibrary{positioning, arrows.meta, shadows.blur, shapes.geometric}
\usepackage{amsmath, amssymb}
\usepackage{url}       % `\url`s
\usepackage{listings}  % Code listings
\usepackage{graphicx}
\usepackage{amsmath}
\usepackage[T1]{fontenc}
\usepackage{xcolor}
\usepackage{minted}
% \usepackage{adjustbox}


% Definicion de pseudocodigo en espaniol
\lstdefinelanguage{PseudocodigoES}{
    morekeywords={INICIO,FIN,PARA,HACER,SI,ENTONCES,FIN_SI,FIN_PARA,IMPRIMIR},
    sensitive=true,
    morecomment=[l]{//},
    morestring=[b]"
}

% Definicion de pseudocodigo en ingles
\lstdefinelanguage{PseudocodeEN}{
    morekeywords={BEGIN,END,FOR,DO,IF,THEN,END_IF,END_FOR,PRINT},
    sensitive=true,
    morecomment=[l]{//},
    morestring=[b]"
}

\lstset{
    basicstyle=\ttfamily\small,
    keywordstyle=\bfseries,
    columns=fullflexible,
    keepspaces=true,
    showstringspaces=false,
    frame=single,
    xleftmargin=0em,
    framexleftmargin=0em,
    framesep=2pt,
    gobble=4,
    resetmargins=true,
    aboveskip=0pt,
    belowskip=0pt
}

\lstdefinestyle{pythonstyle}{
    language=Python,
    basicstyle=\ttfamily\small,
    keywordstyle=\bfseries,
    commentstyle=\itshape,
    stringstyle=\ttfamily,
    showstringspaces=false,
    columns=fullflexible,
    keepspaces=true,
    frame=single,
    breaklines=true,
    xleftmargin=0em,
    framexleftmargin=0em,
    framesep=2pt,
    aboveskip=0pt,
    belowskip=0pt,
    gobble=4,
    resetmargins=true,
    emph={print,input,int,float,str,bool,len,range},
    emphstyle=\bfseries
}

\lstset{style=pythonstyle}

% Estilos del diagrama
\tikzstyle{iniciofin} = [
    ellipse,
    minimum width=2.5cm,
    minimum height=1cm,
    text centered,
    draw=black,
    fill=blue!25
]

\tikzstyle{proceso} = [
    rectangle,
    rounded corners,
    minimum width=3cm,
    minimum height=1cm,
    text centered,
    draw=black,
    fill=cyan!20
]

\tikzstyle{decision} = [
    diamond,
    aspect=2,
    text centered,
    draw=black,
    fill=yellow!35
]

\tikzstyle{flecha} = [
    thick,
    ->,
    >=Stealth
]

% \definecolor{green_python}{HTML}{52A736}

% \newcolumntype{Y}{>{\raggedright\arraybackslash}X}

\frenchspacing
\begin{document}
  \shorthandoff{-}
  \frame[c]{\maketitle}

% \begin{darkframes}

    \begin{frame}{Contenido}
        \begin{itemize}
            \item Introducción
            \item \mintinline{python}{range()}
            \item El ciclo \mintinline{python}{for}
            \item El ciclo \mintinline{python}{while}
            \item \mintinline{python}{for} vs \mintinline{python}{while}
        \end{itemize}
    \end{frame}

    \begin{frame}{Qué es un ciclo (loop)?}
        \begin{itemize}
            \item Un ciclo (loop) es una estructura que permite repetir instrucciones automáticamente
            \item Evita escribir el mismo código muchas veces
            \item Evita tener que realizar la repetición manual
            \item Todo ciclo o bucle, sin excepción, responde a estas dos preguntas:
            \begin{itemize}
                \item ¿Qué se repite? $\rightarrow$ el bloque de instrucciones dentro del bucle
                \item ¿Cuándo se detiene? $\rightarrow$ la condición de salida
            \end{itemize}
        \end{itemize}
    \end{frame}

    \begin{frame}[fragile,t]{Qué es una iteración?}
        \begin{itemize}
            \item Una iteración es cada repetición del ciclo
            \item Si un ciclo se ejecuta 5 veces $\rightarrow$ tiene 5 iteraciones
        \end{itemize}
        \begin{minted}[xleftmargin=-2cm, frame=leftline, fontsize=\small]{python}
            for i in range(3):
                print("Saludo")
        \end{minted}
        \begin{itemize}
            \item Iteración 1 $\rightarrow$ imprime
            \item Iteración 2 $\rightarrow$ imprime
            \item Iteración 3 $\rightarrow$ imprime
        \end{itemize}
    \end{frame}

    \begin{frame}{Para qué se usan los ciclos?}
        \begin{itemize}
            \item Repetir tareas:
            \begin{itemize}
                \item Conteo de soldados
                \item Simulación de rondas de vigilancia
                \item Procesos paso a paso
            \end{itemize}
            \item Automatización
            \item Control de procesos
        \end{itemize}
    \end{frame}

    \begin{frame}{Las tres herramientas}
        \centering
        \begin{tabular}{|c|c|}
            \hline
            Herramienta & ¿Qué hace? \\
            \hline
            \mintinline{python}{for} & Recorre una colección elemento por elemento \\
            \hline
            \mintinline{python}{while} & Repite mientras una condición sea verdadera \\
            \hline
            \mintinline{python}{range} & Genera secuencias de números para iterar \\
            \hline
        \end{tabular}
    \end{frame}

    \begin{frame}[fragile,t]{Tipos de ciclos en Python}
        \begin{itemize}
            \item \mintinline{python}{for} $\rightarrow$ cuando sabemos cuántas veces repetir
            \item \mintinline{python}{while} $\rightarrow$ cuando depende de una condición
        \end{itemize}
    \end{frame}

    \begin{frame}{\mintinline{python}{range()}}
        \begin{itemize}
            \item No es un bucle, es un generador de secuencias de \textbf{números enteros}
            \item No crea una lista en memoria
            \item Es un objeto perezoso (lazy object). No crea todos los números de golpe, los produce uno a uno conforme el ciclo los va pidiendo
        \end{itemize}
    \end{frame}

    \begin{frame}[fragile,t]{\mintinline{python}{range()}}
        \centering
        \begin{tabular}{|c|c|c|}
            \hline
            Caso de uso & Ejemplo código & Salida \\
            \hline
            Básico & \mintinline{python}{range(6)} & \mintinline{python}{0, 1, 2, 3, 4, 5} \\
            \hline
            Start/Stop & \mintinline{python}{range(2, 6)} & \mintinline{python}{2, 3, 4, 5} \\
            \hline
            Step & \mintinline{python}{range(0, 6, 2)} & \mintinline{python}{0, 2, 4} \\
            \hline
            Negative Step & \mintinline{python}{range(6, 0, -1)} & \mintinline{python}{6, 5, 4, 3, 2, 1} \\
            \hline
        \end{tabular}
        
        \begin{minted}[xleftmargin=-2cm, frame=leftline, fontsize=\small]{python}
            range(stop)        # empieza en 0, termina antes de stop
            range(start, stop) # empieza en start, termina antes de stop
            range(start, stop, step) # similar, pero salta de step en step 
        \end{minted}
    \end{frame}

    \begin{frame}[fragile,t]{\mintinline{python}{range()}}
        \begin{itemize}
            \item La notación matemática para el \mintinline{python}{range()} es: \(\text{range}(a,b,c)=\{a+k\cdot c\mid k\in \mathbb{N}_{0}\text{\ y\ }a+k\cdot c\in \text{intervalo\ }[a,b)\}\)
            \item Donde los parámetros corresponden a:
            \begin{itemize}
                \item \(a\) (start): Valor inicial (por defecto 0 si solo se da un parámetro).
                \item \(b\) (stop): Límite de parada (exclusivo, es decir, abierto por la derecha \([a, b)\)).
                \item \(c\) (step): Incremento o paso (por defecto 1).
                \item \(k\): Índice entero de la iteración que comienza en 0 (\(k = 0, 1, 2, \dots\)).
            \end{itemize}
        \end{itemize}
        \begin{minted}[xleftmargin=-2cm, frame=leftline, fontsize=\small]{python}
            range(5)  →  0, 1, 2, 3, 4   # el 5 NO aparece 
        \end{minted}
        \begin{minted}[xleftmargin=-2cm, frame=leftline, fontsize=\small]{python}
            # Generar secuencias del 1 al 10
            range(10):      # ← llega solo hasta el 9
        \end{minted}
        \begin{minted}[xleftmargin=-2cm, frame=leftline, fontsize=\small]{python}
            # Generar secuencias del 1 al 10
            range(1, 11):      # ← correcto
        \end{minted}
    \end{frame}

    \begin{frame}[fragile,t]{Actividad 1: completar el código}
        \begin{minted}[xleftmargin=-2cm, frame=leftline, fontsize=\small]{python}
            # 1. Generar secuencias del 0 al 8
            range(_____)  # ← 0, 1, 2, 3, 4, 5, 6, 7, 8
            # 2. Generar secuencias del 4 al 8
            range(_____)  # ← 4, 5, 6, 7, 8
            # 3. Generar secuencias del 2 al 8, con salto de 2
            range(_____)  # ← 2, 4, 6, 8
            # 4. Generar secuencias del 8 al 2, en orden decreciente
            range(_____)  # ← 8, 7, 6, 5, 4, 3, 2
            # 5. Generar secuencias del 8 al 2, con salto de 2
            range(_____)  # ← 8, 6, 4, 2
        \end{minted}
    \end{frame}

    \begin{frame}{Ciclo \mintinline{python}{for}}
        \begin{itemize}
            \item Itera sobre una secuencias u otros objetos iterables
            \item El código a repetir debe estar indentado
            \item Recorre cada elemento, uno por uno y ejecuta el bloque de código con él
            \item Maneja el inicio, el conteo y el final
        \end{itemize}
    \end{frame}

    \begin{frame}[fragile,t]{Estructura de \mintinline{python}{for}}
        \begin{minted}[xleftmargin=-2cm, frame=leftline, fontsize=\small]{python}
            for elemento in secuencia:
                instrucciones
        \end{minted}

        \begin{minted}[xleftmargin=-2cm, frame=leftline, fontsize=\small]{python}
            for variable in range(...):
                instrucciones
        \end{minted}
    \end{frame}

    \begin{frame}[fragile,t]{Ejemplos usando \mintinline{python}{for}}
        \begin{minted}[xleftmargin=-2cm, frame=leftline, fontsize=\small]{python}
            for i in range(7):
                print(i)   # ← 0, 1, 2, 3, 4, 5, 6
        \end{minted}
        \begin{minted}[xleftmargin=-2cm, frame=leftline, fontsize=\small]{python}
            for i in range(2, 10):
                print(i)   # ← 2, 3, 4, 5, 6, 7, 8, 9
        \end{minted}
        \begin{minted}[xleftmargin=-2cm, frame=leftline, fontsize=\small]{python}
            for i in range(10, 1, -2):
                print(i)   # ← 10, 8, 6, 4, 2
        \end{minted}
        \begin{minted}[xleftmargin=-2cm, frame=leftline, fontsize=\small]{python}
            for i in "hello":
                print(i)   # ← h, e, l, l, o
        \end{minted}
    \end{frame}

    \begin{frame}[fragile,t]{Ejemplos usando \mintinline{python}{for}}
        Contar estudiantes
        \begin{minted}[xleftmargin=-2cm, frame=leftline, fontsize=\small]{python}
            for i in range(1, 6):
                print("Estudiante", i)
        \end{minted}
        Simular entregas de exámenes
        \begin{minted}[xleftmargin=-2cm, frame=leftline, fontsize=\small]{python}
            for i in range(3):
                print("Examen entregado")
        \end{minted}
        Cuenta regresiva
        \begin{minted}[xleftmargin=-2cm, frame=leftline, fontsize=\small]{python}
            for i in range(5, 0, -1):
                print("Cuenta:", i)
            print("Terminó la cuenta regresiva")
        \end{minted}
    \end{frame}

    \begin{frame}[fragile,t]{Actividad 2: Completar el código}
        \begin{minted}[xleftmargin=-2cm, frame=leftline, fontsize=\small]{python}
            # Imprimir 10 veces el texto: Universidad Católica de Chile
            for i in range(____):
                print(________)

            # Sumar los números del 1 al 10 e imprimir el resultado
            suma = 0
            for i in range(_____):
                suma = suma + _________
            print(_________)

            # Sumar los 4 primeros números pares, sin incluir el 0, e imprimir el resultado
            resultado = 0
            for numero_par in range(_____):
                resultado = resultado + _________
            print(__________)
        \end{minted}
    \end{frame}

    \begin{frame}{Aspectos importantes de \mintinline{python}{for}}
        \begin{itemize}
            \item Número de iteraciones conocido
            \item Variable de control cambia automáticamente
            \item En Python, todo lo que sea iterable puede usarse con \mintinline{python}{for}
        \end{itemize}
    \end{frame}

    \begin{frame}{Ciclo \mintinline{python}{while}}
        \begin{itemize}
            \item Repite un bloque de código \textbf{mientras una condición sea verdadera}
            \item No se sabe de antemano cuántas veces va a ejecutarse
            \item El ciclo \mintinline{python}{while} necesita ser controlado
            \item Es necesario crear la condición que se evalúa en el ciclo
            \item La condición tiene que ser modificada dentro del ciclo
            \item \textbf{Si la condición no cambia se genera un ciclo infinito}
        \end{itemize}
    \end{frame}

    \begin{frame}{\mintinline{python}{while}: ¿Cómo funciona?}
        El \mintinline{python}{while} sigue siempre este ciclo:
        \begin{enumerate}
            \item Evalúa la condición $\rightarrow$ ¿es verdadera?
            \item Sí $\rightarrow$ ejecuta las instrucciones $\rightarrow$ vuelve al paso 1
            \item No $\rightarrow$ el ciclo termina
        \end{enumerate}
        condición verdadera $\rightarrow$ ejecutar $\rightarrow$ volver a evaluar

        condición falsa     $\rightarrow$ salir del bucle
    \end{frame}

    \begin{frame}[fragile,t]{Estructura de \mintinline{python}{while}}
        \begin{minted}[xleftmargin=-2cm, frame=leftline, fontsize=\small]{python}
            while condicion:
                instrucciones
        \end{minted}
    \end{frame}

    \begin{frame}[fragile,t]{Ejemplos con \mintinline{python}{while}}
        Contar
        \begin{minted}[xleftmargin=-2cm, frame=leftline, fontsize=\small]{python}
            contador = 0
            while contador < 5:
                print(contador)
                contador += 1
        \end{minted}
        Esperar orden
        \begin{minted}[xleftmargin=-2cm, frame=leftline, fontsize=\small]{python}
            orden = "no"
            while orden != "si":
                print("Esperando orden...")
                orden = input("Ingrese orden (si/no): ")
        \end{minted}
        Nivel de bencina
        \begin{minted}[xleftmargin=-2cm, frame=leftline, fontsize=\small]{python}
            bencina = 100
            while bencina > 0:
                print("Manejando 10 km")
                bencina = bencina - 20
        \end{minted}
    \end{frame}

    \begin{frame}[fragile,t]{Actividad 3: completar el código}
        \begin{minted}[xleftmargin=-2cm, frame=leftline, fontsize=\small]{python}
            # imprimir el contador mientras contador sea mayor a 4
            contador = 10
            while ____________:
                print(contador)
                contador = contador - 1
            
            # imprimir: "password invalido" mientras el password ingresado sea distinto de "passwd123" 
            password = input()
            while _____________:
                print("password invalido")
                password = input("Ingrese el password: ")
        \end{minted}
    \end{frame}

    \begin{frame}{Aspectos importantes de \mintinline{python}{while}}
        \begin{itemize}
            \item Puede ser infinito si no se controla
            \item Depende de una condición
            \item Requiere actualizar la condición o las variables manualmente
            \item \textbf{Debe tener al menos una línea que acerque al programa a la condición de salida}
        \end{itemize}
    \end{frame}

    \begin{frame}[fragile,t]{Ejemplos con \mintinline{python}{while}}
        \begin{minted}[xleftmargin=-2cm, frame=leftline, fontsize=\small]{python}
            # Ciclo infinito — contador nunca cambia
            contador = 0
            while contador < 5:
                print("hola")      # ← contador sigue siendo 0 para siempre
        \end{minted}

        \begin{minted}[xleftmargin=-2cm, frame=leftline, fontsize=\small]{python}
            # Correcto — contador avanza cada vuelta
            contador = 0
            while contador < 5:
                print("hola")
                contador += 1      # ← esto es lo que permite que termine
        \end{minted}
    \end{frame}

    \begin{frame}{Diferencias \mintinline{python}{for} vs \mintinline{python}{while}}
        \centering
        \begin{tabular}{|c|c|c|}
            \hline
             & \mintinline{python}{for} & \mintinline{python}{while} \\
            \hline
            Pregunta qué responde & ¿Qué recorro? & ¿Cuándo paro? \\
            \hline
            Se detiene cuando .. & La colección se agota & La condición es falsa \\
            \hline
            ¿Gestionas el contador? & No - automático & Sí - manual \\
            \hline
            Riesgo principal & Casi ninguno & Bucle infinito \\
            \hline
        \end{tabular}
        \mintinline{python}{for} $\rightarrow$ sabes \textbf{qué} estás recorriendo, \mintinline{python}{while} $\rightarrow$ sabes \textbf{cuándo} quieres parar 
    \end{frame}

    \begin{frame}{Cuándo usar cada uno?}
        \begin{itemize}
            \item Usar \mintinline{python}{for} cuando:
            \begin{itemize}
                \item Sabes cuántas veces repetir
                \item Ej: número de estudiantes, exámenes, pasos
            \end{itemize}
            \item Usar \mintinline{python}{while} cuando:
            \begin{itemize}
                \item No sabes cuántas veces
                \item Depende de una condición
                \item Ej: esperar orden, monitorear estado
            \end{itemize}
        \end{itemize}
    \end{frame}

    \begin{frame}{Mini actividad}
        \begin{itemize}
            \item Repetir 10 veces el saludo '¡Hola, estudiante N° ...! $\rightarrow$ Usa \mintinline{python}{for}
            \item Valida la contraseña ingresada; la contraseña tiene que ser igual a 'puc1234' $\rightarrow$. Usa \mintinline{python}{while}
            \item Cuenta regresiva de lanzamiento: Desde 10 hasta 0 $\rightarrow$ Usar \mintinline{python}{for}
        \end{itemize}
    \end{frame}

    \begin{frame}[fragile,t]{Resumen}
        \begin{itemize}
            \item Un ciclo permite automatizar repeticiones
            \item \mintinline{python}{for} $\rightarrow$ controlado
            \item \mintinline{python}{while} $\rightarrow$ condicionado
            \item Ambos son fundamentales para construir algoritmos
        \end{itemize}
    \end{frame}
\end{document}
